\documentclass[11pt,a4paper]{article}

\usepackage[T1]{fontenc}
\usepackage[utf8]{inputenc}
\usepackage{lmodern}

\usepackage[a4paper,margin=1in]{geometry}
\usepackage{setspace}
\setstretch{1.12}

\setlength{\parindent}{15pt}
\setlength{\parskip}{0pt}

\usepackage{lmodern}
\usepackage{amsmath,amssymb,amsfonts}
\usepackage{booktabs}
\usepackage{longtable}
\usepackage{array}
\usepackage{tabularx}
\usepackage{enumitem}
\usepackage{hyperref}
\usepackage{setspace}
\usepackage{ragged2e}

\hypersetup{
	colorlinks=true,
	linkcolor=black,
	citecolor=black,
	urlcolor=blue,
	pdftitle={Acyclicity Statement and Dependency Criteria for the GTCS},
	pdfauthor={Kostiantyn Osmolovskyi}
}

\setlist[itemize]{topsep=3pt,itemsep=2pt,parsep=0pt}
\setlist[enumerate]{topsep=3pt,itemsep=2pt,parsep=0pt}


\title{Acyclicity Statement and Dependency Criteria\\ for the General Theory of Cognitive Structuring}
	
\author{
	Kostiantyn Osmolovskyi\\
	\small Independent Researcher, Ukraine\\
	\small ORCID: 0009-0006-3144-7237\\
	\small \texttt{constantinosmol@gmail.com}
}

\date{}

% ---------------------------------------------------------------------------
% Part of a series: General Theory of Cognitive Structuring (GTCS)
% Autor: Kostiantyn Osmolovskyi (ORCID: 0009-0006-3144-7237)
% License: Creative Commons Attribution 4.0 International (CC BY 4.0)
% Repository: https://github.com/k-osmolovskyi/k-osmolovskyi.github.io
% DOI: https://doi.org/10.5281/zenodo.19701825
% ---------------------------------------------------------------------------

\begin{document}
	\maketitle

\begin{center}
	\small
	Technical Note No. 2026-3\\
	Version 1.0
\end{center}

\vspace{0.5em}
	
	\maketitle
	
	\section{Purpose of this note}
	
	This note provides a public dependency statement for the paper series associated with the
	\emph{General Theory of Cognitive Structuring}. Its purpose is not to restate the formal
	content of the individual papers, but to clarify the dependency criteria under which the
	series may be read as a non-circular theoretical development.
	
	More specifically, the note is intended to support external verification by:
	\begin{itemize}[leftmargin=2em]
		\item distinguishing logical dependency from textual cross-reference,
		\item specifying what counts as a dependency node,
		\item specifying what counts as a dependency edge,
		\item clarifying in what sense the series is acyclic,
		\item identifying common sources of false cycle attribution.
	\end{itemize}
	
	The note is public and self-contained at the level of dependency criteria. It is not intended
	to replace the papers themselves, nor to function as a full dependency registry.
	
	\section{Scope and limits}
	
	This document is limited to the dependency logic of the series.
	
	It does \emph{not} attempt to provide:
	\begin{itemize}[leftmargin=2em]
		\item a full glossary of terms,
		\item a full node-by-node registry,
		\item a complete theorem index,
		\item a proof compendium,
		\item a bibliographic reconstruction of every cross-reference.
	\end{itemize}
	
	Its narrower aim is to specify the criteria under which the core definitional, operator-level,
	construction-level, and proposition-level structure of the series may be treated as a directed
	acyclic graph.
	
	Detailed node registries, paper-level dependency tables, and extended graph materials may be
	provided separately in a supporting dependency companion.
	
	\section{Acyclicity Statement v1}
	
	\subsection{Core claim}
	
	Under the dependency criteria adopted in this note, the core definitional and proposition-level
	structure of the \emph{General Theory of Cognitive Structuring} series is representable as a
	directed acyclic graph.
	
	This claim is intentionally narrower than the claim that the series contains no cross-reference,
	no retrospective clarification, or no synthetic reformulation. It concerns only \emph{logical
		dependency} in the sense defined below.
	
	\subsection{What counts as acyclicity}
	
	Acyclicity is evaluated only at the level of logical dependency.
	
	A cycle is present only if there exists a sequence of nodes
	\[
	N_1 \rightarrow N_2 \rightarrow \dots \rightarrow N_k \rightarrow N_1
	\]
	such that each arrow expresses a genuine logical dependency.
	
	Accordingly, the following do \emph{not} by themselves constitute a cycle:
	\begin{itemize}[leftmargin=2em]
		\item mutual textual references,
		\item retrospective clarification of earlier operational usage,
		\item later formal grounding of notions previously used in a limited way,
		\item synthetic integration of earlier branches,
		\item bibliographic cross-reference for orientation or context.
	\end{itemize}
	
	\subsection{Dependency criterion}
	
	For the purposes of this note, node $A$ depends logically on node $B$ only if at least one of
	the following conditions holds:
	
	\begin{enumerate}[leftmargin=2em]
		\item \textbf{Definitional dependence.}  
		$A$ cannot be stated without invoking $B$ as part of its definition.
		
		\item \textbf{Operator dependence.}  
		$A$ is an operator, mapping, or gate whose domain, codomain, or arguments require $B$.
		
		\item \textbf{Construction dependence.}  
		$A$ is a construction obtained from $B$ or from a set including $B$.
		
		\item \textbf{Proposition/theorem dependence.}  
		$A$ is a proposition, theorem, or formal consequence whose statement or proof requires $B$.
		
		\item \textbf{Architectural-layer dependence.}  
		$A$ is introduced explicitly as operating on the output, domain, or admissible image generated by $B$.
	\end{enumerate}
	
	\subsection{What does not count as dependency}
	
	The following do not count as logical dependency for the purposes of acyclicity:
	
	\subsubsection*{Textual cross-reference}
	
	A reference to another paper for context, positioning, terminology, or overview does not create
	a logical edge.
	
	\subsubsection*{Retrospective clarification}
	
	A later paper may formally clarify the architectural source of a notion that earlier papers used
	operationally. This does not automatically mean that the earlier paper logically depended on the
	later one.
	
	\subsubsection*{Synthetic integration}
	
	A synthetic paper that gathers earlier branches into a single framework is treated as a convergence
	node, not as the foundational source of the branches it integrates.
	
	\subsubsection*{Broad thematic relevance}
	
	Thematic relevance or conceptual proximity is not sufficient for a dependency edge unless one of
	the dependency criteria stated above is satisfied.
	
	\subsection{Provisional conclusion}
	
	The present corpus supports the following v1 conclusion:
	
	\begin{quote}
		\textbf{Acyclicity Statement v1.} Under the dependency criteria adopted in this document, the core definitional, operator-level, construction-level, and proposition-level structure of the current \emph{General Theory of Cognitive Structuring} series is representable as a directed acyclic graph. Apparent cycles produced by retrospective clarification, synthetic integration,	or textual cross-reference do not count as logical cycles.
	\end{quote}
	
	This is a provisional v1 statement because some node boundaries may still be refined in future
	supporting materials. However, the current paper sequence already supports the absence of core
	circularity at the level defined here.
	
	\section{Why apparent cycles do not count as cycles}
	
	\subsection{Coherence and invariants}
	
	An apparent cycle may be inferred when an earlier paper defines coherence relative to a stability
	region and a later paper provides the stronger architectural grounding of that stability region
	through invariant structure.
	
	Under the present criteria, this does not count as a logical cycle. The earlier paper introduces
	a stability-region-based geometric definition. The later paper refines the architectural grounding
	of that region. Refinement is not treated as retroactive logical generation.
	
	\subsection{Coherence representation and pre-symbolic admissibility}
	
	An apparent cycle may also be inferred when one paper introduces coherence representation as an
	internal regulatory variable and a later paper states that representation is constrained by a
	pre-symbolic admissibility layer.
	
	Under the present criteria, this does not count as a cycle. The later paper narrows the upstream
	domain through which representation may be constructed. It does not become the foundational source
	of the earlier representational definition.
	
	\subsection{Restricted accessibility and coherence}
	
	Restricted accessibility is a consequence drawn from geometric coherence together with admissibility-constrained
	regulatory access. The fact that it later clarifies how coherence should be interpreted does not make the
	original geometric notion of coherence logically dependent on the restricted-accessibility result.
	
	\subsection{Synthetic theory and earlier papers}
	
	A synthetic theory paper may gather earlier branches into a unified layered framework. This does
	not create reverse foundational dependency from the earlier papers to the synthetic integration.
	Otherwise any mature synthetic paper would generate a false cycle with all earlier work it organizes.
	
	\section{Minimal unfolding of the series}
	
	The present note does not include a full dependency table. However, the paper series may be read
	through the following minimal unfolding order:
	
	\begin{enumerate}[leftmargin=2em]
		\item entry distinction between processing and structural updating;
		\item geometric formalization of coherence;
		\item formalization of structural admissibility in the joint regulatory state;
		\item formalization of overload and trajectory dependence;
		\item formalization of invariants as architectural constraints;
		\item formalization of compression as admissible structural transformation;
		\item introduction of coherence representation as an internal regulatory variable;
		\item introduction of pre-symbolic admissibility as an earlier filtering layer;
		\item derivation of restricted accessibility of coherence;
		\item extension to multi-system regulation and inter-system conflict;
		\item synthetic integration in the \emph{General Theory of Cognitive Structuring}.
	\end{enumerate}
	
	This ordering is sufficient for the present note because it shows that later layers are introduced
	as extensions, refinements, or integrations of earlier layers rather than as their hidden logical
	preconditions.
	
	\section{Node and Edge Taxonomy}
	
	\subsection{Purpose}
	
	This taxonomy defines the admissible node types and edge types used in the dependency note.
	
	Its purpose is to keep the graph readable and to prevent category mistakes, especially confusion
	between:
	\begin{itemize}[leftmargin=2em]
		\item definitions and constructions,
		\item operators and propositions,
		\item architectural objects and representational objects,
		\item synthetic organizing nodes and theorem-grade nodes.
	\end{itemize}
	
	\subsection{Node taxonomy}
	
	\subsubsection*{D --- Definition}
	
	A node of type \textbf{D} introduces a concept whose role is primarily definitional.
	
	Use type D when the node answers: \emph{What is this?}
	
	\subsubsection*{O --- Operator}
	
	A node of type \textbf{O} is a map, functional, gate, or formal operator acting on a specified domain.
	
	Use type O when the node's identity depends on input-output role or formal action.
	
	\subsubsection*{C --- Construction}
	
	A node of type \textbf{C} is a derived object, region, domain, state variable, or structured set
	generated from prior nodes.
	
	Use type C when the node is a formal object built within the framework but is neither a basic
	definition nor a result statement.
	
	\subsubsection*{P --- Proposition / proposition-level result}
	
	A node of type \textbf{P} states a derived consequence, lemma-like claim, or proposition-level relation.
	
	Use type P when the node answers: \emph{What follows?}
	
	\subsubsection*{T --- Theorem}
	
	A node of type \textbf{T} is a formally marked theorem or theorem-like result whose proof role
	matters independently in the graph.
	
	Use type T sparingly.
	
	\subsubsection*{R --- Representational variable or representational construction}
	
	A node of type \textbf{R} is reserved for internal variables whose role is specifically representational
	within regulation.
	
	Use type R when representational status itself is central to the node.
	
	\subsubsection*{A --- Assumption / axiom-like condition}
	
	A node of type \textbf{A} represents an assumption, regularity condition, monotonicity condition,
	or axiom supporting later results.
	
	Use type A when the node is a supporting condition rather than a concept or derived consequence.
	
	\subsubsection*{M --- Meta-structural or graph-organizing node}
	
	A node of type \textbf{M} is used for graph organization at a higher level without forcing a false
	theorem-object reading.
	
	Use type M only when the node is helpful for DAG readability but should not be mistaken for a
	basic formal unit.
	
	\subsection{Edge taxonomy}
	
	\subsubsection*{defines}
	
	Use \textbf{defines} when node $A$ provides the explicit definition of node $B$.
	
	Without the source node, the target node's definition is incomplete.
	
	\subsubsection*{uses}
	
	Use \textbf{uses} when node $B$ invokes node $A$ as part of its statement or construction, but
	$A$ does not uniquely determine $B$.
	
	This is weaker than direct definitional generation.
	
	\subsubsection*{requires}
	
	Use \textbf{requires} when the target node cannot be stated in admissible form without the source node.
	
	This is stronger than \emph{uses} and is close to prerequisite status.
	
	\subsubsection*{refines}
	
	Use \textbf{refines} when a later node makes an earlier node more explicit, narrower, or more
	formally grounded without overturning it.
	
	A \emph{refines} edge does not imply reverse foundational dependence.
	
	\subsubsection*{extends}
	
	Use \textbf{extends} when a node carries an earlier object into a new layer, domain, or setting.
	
	This marks forward expansion rather than basic definition.
	
	\subsubsection*{derives}
	
	Use \textbf{derives} when a proposition or theorem follows from earlier nodes.
	
	This is the preferred edge for result-level dependence.
	
	\subsubsection*{integrates}
	
	Use \textbf{integrates} only for synthetic nodes that gather earlier branches into a convergence
	structure.
	
	An \emph{integrates} edge is non-foundational: it marks convergence, not source-generation.
	
	\section{Edge restrictions for the public note}
	
	To preserve acyclicity and readability, the following restrictions are recommended:
	
	\begin{enumerate}[leftmargin=2em]
		\item A bibliographic citation alone does not warrant an edge.
		\item Synthetic nodes may receive \emph{integrates} edges from earlier nodes, but should not
		generate foundational reverse edges.
		\item Rhetorical preview does not create dependency.
		\item When in doubt, include only the smallest set of direct dependencies needed for correct statement.
		\item A later formal grounding of an earlier operational term should usually be marked as
		\emph{refines}, not \emph{requires}.
	\end{enumerate}
	
	\section{Recommended use of this note}
	
	This document is intended to function as a public, self-contained dependency statement supporting
	external verification of the series.
	
	It should be used:
	\begin{itemize}[leftmargin=2em]
		\item together with the individual papers,
		\item prior to any full node registry or extended dependency dossier,
		\item as a criterion note rather than as a replacement for formal content.
	\end{itemize}
	
	A separate companion document may provide:
	\begin{itemize}[leftmargin=2em]
		\item a full node registry,
		\item a paper-level dependency table,
		\item extended graph annotations,
		\item and appendix-level audit materials.
	\end{itemize}
	
	\section{Compact appendix form}
	
	\subsection*{Node types}
	
	\begin{itemize}[leftmargin=2em]
		\item \textbf{D} --- Definition
		\item \textbf{O} --- Operator
		\item \textbf{C} --- Construction
		\item \textbf{P} --- Proposition
		\item \textbf{T} --- Theorem
		\item \textbf{R} --- Representational variable / construction
		\item \textbf{A} --- Assumption / axiom-like condition
		\item \textbf{M} --- Meta-structural / graph-organizing node
	\end{itemize}
	
	\subsection*{Edge types}
	
	\begin{itemize}[leftmargin=2em]
		\item \textbf{defines} --- source gives the target's definition
		\item \textbf{requires} --- target cannot be stated without source
		\item \textbf{uses} --- target functionally invokes source
		\item \textbf{refines} --- target narrows or grounds source more explicitly
		\item \textbf{extends} --- target carries source into a new layer or domain
		\item \textbf{derives} --- target follows as a proposition/theorem consequence
		\item \textbf{integrates} --- synthetic convergence relation only
	\end{itemize}
	
\end{document}
